\documentclass[12pt, a4paper]{article}

% ── Encodage & langue ──────────────────────────────────────────────────────────
\usepackage[utf8]{inputenc}
\usepackage[T1]{fontenc}
\usepackage[french]{babel}

% ── Mise en page ───────────────────────────────────────────────────────────────
\usepackage[top=2.5cm, bottom=2.5cm, left=2.5cm, right=2.5cm]{geometry}
\usepackage{setspace}
\usepackage{pdflscape}   % page en paysage
\onehalfspacing

% ── Typographie ────────────────────────────────────────────────────────────────
\usepackage{lmodern}
\usepackage{microtype}

% ── Couleurs & graphiques ──────────────────────────────────────────────────────
\usepackage{xcolor}
\usepackage{graphicx}
\usepackage{tikz}

\definecolor{primary}{RGB}{30, 80, 160}
\definecolor{lightgray}{RGB}{245, 245, 245}
\definecolor{axeA}{RGB}{30,  80, 160}
\definecolor{axeB}{RGB}{0,  140, 140}
\definecolor{axeC}{RGB}{0,  140, 100}

% ── En-têtes & pieds de page ───────────────────────────────────────────────────
\usepackage{fancyhdr}
\pagestyle{fancy}
\fancyhf{}
\fancyhead[L]{\small\textcolor{primary}{\textit{Cahier des charges}}}
\fancyhead[R]{\small\textcolor{gray}{\thepage}}
\fancyfoot[C]{\small\textcolor{gray}{\today}}
\renewcommand{\headrulewidth}{0.4pt}
\renewcommand{\headrule}{\hbox to\headwidth{\color{primary}\leaders\hrule height \headrulewidth\hfill}}

% ── Titres de sections ─────────────────────────────────────────────────────────
\usepackage{titlesec}
\titleformat{\section}
  {\color{primary}\large\bfseries}
  {\thesection.}{0.6em}{}
  [\color{primary}\titlerule]
\titleformat{\subsection}
  {\color{primary!80!black}\normalsize\bfseries}
  {\thesubsection.}{0.5em}{}
\titleformat{\subsubsection}
  {\normalsize\bfseries}
  {\thesubsubsection.}{0.5em}{}

% ── Table des matières ─────────────────────────────────────────────────────────
\usepackage[hidelinks]{hyperref}
\usepackage{tocloft}
\renewcommand{\cftsecfont}{\bfseries\color{primary}}
\renewcommand{\cftsecpagefont}{\bfseries\color{primary}}

% ── Tableaux ───────────────────────────────────────────────────────────────────
\usepackage{booktabs}
\usepackage{tabularx}
\usepackage{array}
\usepackage{multirow}
\usepackage{colortbl}
\newcolumntype{L}[1]{>{\raggedright\arraybackslash}p{#1}}
\newcolumntype{C}[1]{>{\centering\arraybackslash}p{#1}}

% ── Listes ─────────────────────────────────────────────────────────────────────
\usepackage{enumitem}
\setlist[itemize]{leftmargin=1.5em, itemsep=0.3em}
\setlist[enumerate]{leftmargin=1.5em, itemsep=0.3em}

% ── Boîtes info ────────────────────────────────────────────────────────────────
\usepackage{tcolorbox}
\tcbuselibrary{skins, breakable}

\newtcolorbox{infobox}[1][]{
  colback=lightgray,
  colframe=primary,
  fonttitle=\bfseries,
  title=#1,
  breakable,
  arc=4pt,
  boxrule=0.6pt,
}

\newtcolorbox{axebox}[2][primary]{
  colback=#1!8,
  colframe=#1,
  fonttitle=\bfseries\color{white},
  title=#2,
  coltitle=white,
  breakable,
  arc=4pt,
  boxrule=0.8pt,
}

% ── Diagramme de Gantt ─────────────────────────────────────────────────────────
\usepackage{pgfgantt}
\usepackage{amssymb}

% ══════════════════════════════════════════════════════════════════════════════
%  PAGE DE TITRE
% ══════════════════════════════════════════════════════════════════════════════
\begin{document}

\begin{titlepage}
  \centering
  \vspace*{1.5cm}

  {\color{primary}\rule{\linewidth}{1.5pt}}\vspace{0.4cm}
  {\color{primary}\rule{\linewidth}{0.4pt}}

  \vspace{1cm}
  {\Huge\bfseries\color{primary} Cahier des charges}\\[0.4cm]
  {\Large\color{gray} Projet fil rouge - Travail de recherche}

  \vspace{0.8cm}
  {\color{primary}\rule{\linewidth}{0.4pt}}\vspace{0.2cm}
  {\color{primary}\rule{\linewidth}{1.5pt}}

  \vspace{2cm}

  \begin{tcolorbox}[
    colback=primary!10, colframe=primary,
    arc=6pt, boxrule=1pt,
    width=0.9\linewidth
  ]
    \centering
    {\large\bfseries\color{primary} Génération automatique de plans de tests}\\[0.3em]
    {\large\bfseries\color{primary} depuis des normes techniques grâce aux LLMs}
  \end{tcolorbox}

  \vspace{2.5cm}

  \begin{tabular}{rl}
    \textbf{Auteur}        & Luca Ceccarelli \\[0.4em]
    \textbf{Formation}     & INFRES 17 \\[0.4em]
    \textbf{Encadrants}    & Christelle Urtado \& Sylvain Vauttier \\[0.4em]
    \textbf{Date}          & Juin 2026 \\[0.4em]
    \textbf{Version}       & 1.0 \\
  \end{tabular}

\end{titlepage}

% ══════════════════════════════════════════════════════════════════════════════
\tableofcontents
\newpage

% ══════════════════════════════════════════════════════════════════════════════
\section{Contexte et périmètre du projet}
% ══════════════════════════════════════════════════════════════════════════════

\subsection{Contexte général}

Dans de nombreux secteurs industriels (télécommunications, finance, santé, automobile),
les produits logiciels doivent se conformer à des normes techniques internationales telles
que les RFC de l'IETF ou les normes ISO/IEC. Un produit peut être développé en suivant
une telle norme, et un \textbf{plan de test} permet de vérifier que ce produit respecte
correctement l'ensemble des spécifications, internes et externes, qui y sont définies.

La norme \textbf{ISO/IEC/IEEE 29119-3:2021} définit le format de référence pour les plans
de test : chaque cas de test y spécifie un identifiant, un objectif, des préconditions, un
scénario et des résultats attendus, sans produire de code exécutable.

\subsection{Enjeu industriel et motivation}

Aujourd'hui, les équipes de conformité passent des semaines à rédiger manuellement des
plans de test à partir de normes qui peuvent compter plusieurs centaines de pages. Ce
travail est coûteux, répétitif et source d'oublis : de nombreuses normes ne disposent ainsi
d'aucun plan de test, ce qui nuit à l'interopérabilité des systèmes par conception.

L'automatisation de ce processus permettrait de réduire significativement le
\textit{time-to-compliance}, de limiter les oublis et d'améliorer l'interopérabilité entre
implémentations. La norme \textbf{ISO/IEC 18013-5} (permis de conduire numérique mDL)
constitue un exemple concret : elle dispose d'un plan de test existant, référencé dans
l'ISO/IEC 18013-6, ce qui offre une base de comparaison idéale pour évaluer une approche
automatisée.

% ══════════════════════════════════════════════════════════════════════════════
\section{Problématique}
% ══════════════════════════════════════════════════════════════════════════════

\subsection{Formulation de la problématique}

\begin{infobox}[Problématique centrale]
\textit{Comment formaliser une approche hybride, combinant l'extraction d'exigences par
NLP et une architecture agentique, pour automatiser la génération de plans de tests
fiables à partir de normes techniques ?}
\end{infobox}

\subsection{Limites des solutions existantes}

Une revue bibliographique systématique (6 articles retenus sur 10 identifiés, issus de
Google Scholar, IEEE Xplore, ACM DL et HAL, mai 2026) a permis d'identifier deux familles
d'approches existantes et leurs limites respectives.

\begin{itemize}
  \item \textbf{Approches LLM existantes} (Haldar \& Capretz 2026 ; Hasan et al. 2025 ;
    Xue et al. 2024 ; Li et al. 2025) : elles génèrent des tests à partir de code source,
    de dépôts GitHub, ou de documents de spécification projet en texte libre. Aucune ne
    part d'une norme technique dense en amont du développement.

  \item \textbf{Approches NLP classiques} (Wang et al. 2022 ; Dwarakanath \& Sengupta
    2012) : elles traitent des cas d'utilisation UML ou des spécifications fonctionnelles
    simples. Elles sont incapables d'absorber la densité et la structure hiérarchique
    d'une norme normative (références croisées, clauses SHALL/SHOULD, annexes normatives).
\end{itemize}

L'angle mort identifié est donc le suivant : \textbf{aucune approche ne génère
un plan de test au format ISO 29119-3 directement depuis une norme technique en amont
de tout développement}, en prenant en compte les dépendances inter-normes.

% ══════════════════════════════════════════════════════════════════════════════
\section{Solution proposée et axes de recherche}
% ══════════════════════════════════════════════════════════════════════════════

\subsection{Vue d'ensemble de l'approche}

La solution proposée repose sur un pipeline en trois axes de recherche successifs et
prioritisés, appliqué initialement à la norme ISO/IEC 18013-5 :

\begin{axebox}[axeA]{Axe A — Extraction des exigences techniques - Priorité 1}
Transformer une norme technique (RFC, ISO) en base de connaissances exploitable, avec
une perte d'information minimale.

\medskip
\textbf{Briques technologiques envisagées :}
\begin{itemize}
  \item NLP déterministe : tokenisation, POS tagging, extraction de règles normatives ;
  \item Chunking sémantique par section normative (LangChain RecursiveCharacterTextSplitter) ;
  \item Embedding et indexation vectorielle pour le retrieval (FAISS).
\end{itemize}

\medskip
\textbf{Plan d'expérience :} test de complétude (aucune exigence perdue), métriques
Précision / Rappel / F1 sur une référence annotée manuellement ; comparaison chunking
naïf vs.\ chunking sémantique.
\end{axebox}

\vspace{0.8em}

\begin{axebox}[axeB]{Axe B — Génération agentique du plan de test - Priorité 2}
Système multi-agents raisonnant sur les exigences extraites (Axe A) pour produire un
plan de test conforme à l'ISO/IEC/IEEE 29119-3:2021.

\medskip
\textbf{Architecture multi-agents :}
\begin{itemize}
  \item \textbf{Agent Analyste} : reçoit une exigence isolée (sortie Axe A), la contextualise ;
  \item \textbf{Agent Documentaliste (RAG)} : interroge la base vectorielle et les sources
    en ligne (RFC, normes référencées) ;
  \item \textbf{Agent Concepteur} : rédige le cas de test au format ISO 29119-3:2021
    (ID, objectif, préconditions, étapes, résultat attendu).
\end{itemize}

\medskip
\textbf{Plan d'expérience :} test de sensibilité entre différents LLMs, évaluation
humaine (pertinence et complétude), couverture ISO (ratio sections couvertes / total),
et validation par comparaison avec le plan de test existant ISO/IEC 18013-6 (En générant le premier test sur la norme ISO/IEC 18013-5).
\end{axebox}

\vspace{0.8em}

\begin{axebox}[axeC]{Axe C — Validation agentique du plan de test - Priorité 3}
Valider le plan généré via un système dédié à la détection d'erreurs, de contradictions
et d'hallucinations.

\medskip
\textbf{Briques technologiques envisagées :} LLM-as-a-Judge, framework Guardrails AI, scoring automatique de conformité au format ISO 29119-3.

\medskip
\textbf{Plan d'expérience :} adversarial testing (injection délibérée d'erreurs),
mesure du taux de détection (contradictions, oublis, hallucinations), comparaison
validation agentique vs.\ relecture humaine.
\end{axebox}

\subsection{Livrables attendus}

\begin{enumerate}
  \item Pipeline d'extraction des exigences depuis une norme PDF (Axe A), évalué sur
    ISO/IEC 18013-5.
  \item Système multi-agents de génération de plans de test au format ISO 29119-3 (Axe B),
    avec évaluation comparative.
  \item Module de validation agentique (Axe C, si le temps le permet).
  \item Généralisation du pipeline sur au moins 3 normes différentes (Phase 3).
  \item Publication scientifique soumise à une conférence ou un journal du domaine.
\end{enumerate}

% ══════════════════════════════════════════════════════════════════════════════
\section{Organisation et rôles}
% ══════════════════════════════════════════════════════════════════════════════

\subsection{Tableau RACI}

Dans le cadre d'un projet de recherche individuel, la matrice RACI structure la
répartition des responsabilités entre l'auteur, ses encadrants et le jury.

\vspace{0.5em}
\begin{center}
\begin{tabularx}{\linewidth}{|>{\raggedright\arraybackslash}X|C{1.6cm}|C{1.9cm}|C{1.9cm}|C{1.6cm}|}
  \hline
  \rowcolor{primary!15}
  \textbf{Tâche / Activité}
    & \textbf{Luca C.}
    & \textbf{C.~Urtado}
    & \textbf{S.~Vauttier}
    & \textbf{Jury} \\
  \hline
  Veille bibliographique et état de l'art        & R, A & C & C & I \\
  \hline
  Conception de la méthode (axes A, B, C)        & R, A & C & C & — \\
  \hline
  Implémentation Axe A (extraction NLP)          & R, A & C & C & — \\
  \hline
  Implémentation Axe B (multi-agents)            & R, A & C & C & — \\
  \hline
  Implémentation Axe C (validation, optionnel)   & R, A & C & C & — \\
  \hline
  Évaluation et analyse des résultats            & R, A & C & C & I \\
  \hline
  Rédaction du rapport / article                 & R, A & C & C & I \\
  \hline
  Validation des jalons intermédiaires           & R    & A & A & I \\
  \hline
  Présentation et soutenance finale              & R, A & C & C & I \\
  \hline
\end{tabularx}
\end{center}
\vspace{0.4em}
{\small \textbf{Légende :} R = Responsable (réalise), A = Autorité (valide), C = Consulté,
I = Informé.}

% ══════════════════════════════════════════════════════════════════════════════
\section{Jalons et planification}
% ══════════════════════════════════════════════════════════════════════════════

\subsection{Phases du projet}

Le projet se déroule en trois phases, couvrant mars 2026 à mars 2027.

\vspace{0.5em}
\begin{center}
\begin{tabularx}{\linewidth}{|C{0.8cm}|L{3.2cm}|>{\raggedright\arraybackslash}X|L{2.5cm}|}
  \hline
  \rowcolor{primary!15}
  \textbf{N°} & \textbf{Phase} & \textbf{Activités principales} & \textbf{Période} \\
  \hline
  \rowcolor{axeA!10}
  1 & \textbf{Exploration}
    & Définition du sujet avec les encadrants — Recherche bibliographique
      (10 articles) — Lecture critique des 6 articles retenus — Feedbacks
      d'affinement du positionnement — Rédaction de la présentation intermédiaire
    & Mars – Juin 2026 \\
  \hline
  \rowcolor{axeB!10}
  2 & \textbf{Prototype}
    & Rédaction de l'introduction (éléments Phase 1) — Implémentation Axe A
      (extraction ISO/IEC 18013-5) — Évaluation Précision/Rappel/F1 —
      Implémentation Axe B (multi-agents) — Évaluation par comparaison avec
      le plan de test existant — Feedbacks encadrants et ajustements
    & Sept. – Déc. 2026 \\
  \hline
  \rowcolor{axeC!10}
  3 & \textbf{Évaluation \& Rédaction}
    & Axe C si temps disponible — Généralisation sur 3 normes différentes —
      Rédaction de la publication — Soutenance finale
    & Janv. – Mars 2027 \\
  \hline
\end{tabularx}
\end{center}

\subsection{Tableau détaillé des jalons}

\vspace{0.5em}
\begin{center}
\begin{tabularx}{\linewidth}{|C{0.6cm}|>{\raggedright\arraybackslash}X|L{2.2cm}|L{2.0cm}|C{1.5cm}|}
  \hline
  \rowcolor{primary!15}
  \textbf{J\#} & \textbf{Jalon / Tâche} & \textbf{Prérequis} & \textbf{Échéance} & \textbf{Charge} \\
  \hline
  J1 & Veille biblio systématique \& état de l'art (10 articles, sélection de 6)
     & —     & Juin 2026    & \textasciitilde{}3 sem. \\
  \hline
  J2 & Positionnement et définition des axes A/B/C
     & J1    & Juin 2026    & \textasciitilde{}1 sem. \\
  \hline
  \rowcolor{orange!15}
  \multicolumn{5}{|l|}{\textit{Livrable intermédiaire 1 : présentation orale — Juin 2026}} \\
  \hline
  J3 & Implémentation pipeline d'extraction Axe A
     & J2    & Sept. 2026    & \textasciitilde{}2 sem. \\
  \hline
  J4 & Annotation de la référence \& évaluation P/R/F1 (Axe A)
     & J3    & Oct. 2026    & \textasciitilde{}1 sem. \\
  \hline
  J5 & Conception architecture multi-agents (Axe B)
     & J3    & Nov. 2026    & \textasciitilde{}3 sem. \\
  \hline
  J6 & Génération d'un premier plan de test ISO 29119-3 (Axe B)
     & J4, J5 & Nov. 2026  & \textasciitilde{}1 sem. \\
  \hline
  J7 & Évaluation Axe B (couverture ISO, comparaison plan existant, éval. humaine)
     & J6    & Déc. 2026    & \textasciitilde{}2 sem. \\
  \hline
  \rowcolor{orange!15}
  \multicolumn{5}{|l|}{\textit{Livrable intermédiaire 2 : réunion rapport d'avancement — Déc. 2026}} \\
  \hline
  J8 & Axe C — Validation agentique (optionnel, si temps disponible)
     & J7    & Fév. 2027    & \textasciitilde{}1 sem. \\
  \hline
  J9 & Généralisation sur 3 normes différentes
     & J7    & Fév. 2027    & \textasciitilde{}1 sem. \\
  \hline
  J10 & Rédaction du rapport final et de la publication
     & J7    & Mars 2027    & \textasciitilde{}2 sem. \\
  \hline
  J11 & Préparation et soutenance finale
     & J10   & Mars 2027    & \textasciitilde{}1 sem. \\
  \hline
  \rowcolor{orange!15}
  \multicolumn{5}{|l|}{\textit{Livrable final : soutenance finale, publication, code source — Mars 2027}} \\
  \hline
\end{tabularx}
\end{center}

\subsection{Diagramme de Gantt prévisionnel}

Le diagramme de Gantt est présenté en page suivante.

% ── PAGE PAYSAGE POUR LE GANTT ────────────────────────────────────────────────
\begin{landscape}
\thispagestyle{empty}

\vspace*{0.5cm}
{\large\bfseries\color{primary} Diagramme de Gantt prévisionnel}
\vspace{0.5cm}

\begin{center}
\begin{ganttchart}[
  hgrid,
  vgrid={*1{dotted}},
  x unit=1.02cm,
  y unit title=0.7cm,
  y unit chart=0.62cm,
  canvas/.style={fill=white, draw=gray!30},
  title/.style={fill=primary!20, draw=primary!60},
  title label font=\small\bfseries\color{primary},
  title height=1,
  bar height=0.5,
  bar top shift=0.25,
  bar/.style={fill=primary!55, draw=primary!80, rounded corners=1pt},
  bar label font=\small,
  bar label node/.style={text width=4.8cm, align=right, font=\small},
  group/.style={fill=primary!80, draw=primary},
  group label font=\small\bfseries,
  milestone/.style={fill=orange!90, draw=orange!70!black, rounded corners=2pt},
  milestone label font=\small\itshape,
  today=4,
  today rule/.style={draw=red!70, thick, dashed},
  today label={\small\color{red!70}\textit{Aujourd'hui}},
  today label font=\small,
  link/.style={-latex, draw=gray!60}
]{1}{12}

  % ── Titres ────────────────────────────────────────────────────────────────
  \gantttitle{Phase 1 — Exploration}{4}
  \gantttitle{Phase 2 — Prototype}{5}
  \gantttitle{Phase 3 — Éval. \& Rédaction}{3} \\

  \gantttitle[title label font=\footnotesize\color{gray}]{Mar}{1}
  \gantttitle[title label font=\footnotesize\color{gray}]{Avr}{1}
  \gantttitle[title label font=\footnotesize\color{gray}]{Mai}{1}
  \gantttitle[title label font=\footnotesize\color{gray}]{Juin}{1}
  \gantttitle[title label font=\footnotesize\color{gray}]{Sep}{1}
  \gantttitle[title label font=\footnotesize\color{gray}]{Oct}{1}
  \gantttitle[title label font=\footnotesize\color{gray}]{Nov}{1}
  \gantttitle[title label font=\footnotesize\color{gray}]{Déc}{1}
  \gantttitle[title label font=\footnotesize\color{gray}]{Jan}{1}
  \gantttitle[title label font=\footnotesize\color{gray}]{Fév}{1}
  \gantttitle[title label font=\footnotesize\color{gray}]{Mar (s)}{1}
  \gantttitle[title label font=\footnotesize\color{gray}]{Mar (f)}{1} \\

  % ── Phase 1 ───────────────────────────────────────────────────────────────
  \ganttbar[bar/.style={fill=axeA!50,draw=axeA}]
    {J1 — Veille bibliographique}{1}{3} \\
  \ganttbar[bar/.style={fill=axeA!50,draw=axeA}]
    {J2 — Positionnement \& axes}{3}{4} \\
  \ganttmilestone[milestone/.style={fill=orange,draw=orange!70!black}]
    {Livrable 1 : présentation orale}{4} \\[0.2em]

  % ── Phase 2 ───────────────────────────────────────────────────────────────
  \ganttbar[bar/.style={fill=axeB!50,draw=axeB}]
    {J3 — Impl. extraction Axe A}{5}{6} \\
  \ganttbar[bar/.style={fill=axeB!50,draw=axeB}]
    {J4 — Annotation \& éval. P/R/F1}{6}{7} \\
  \ganttbar[bar/.style={fill=axeB!50,draw=axeB}]
    {J5 — Architecture multi-agents}{6}{7} \\
  \ganttbar[bar/.style={fill=axeB!50,draw=axeB}]
    {J6 — Génération plan de test (Axe B)}{7}{8} \\
  \ganttbar[bar/.style={fill=axeB!50,draw=axeB}]
    {J7 — Évaluation Axe B}{8}{8} \\
  \ganttmilestone[milestone/.style={fill=orange,draw=orange!70!black}]
    {Livrable 2 : rapport d'avancement}{8} \\[0.2em]

  % ── Phase 3 ───────────────────────────────────────────────────────────────
  \ganttbar[bar/.style={fill=axeC!45,draw=axeC}]
    {J8 — Axe C validation (optionnel)}{9}{10} \\
  \ganttbar[bar/.style={fill=axeC!45,draw=axeC}]
    {J9 — Généralisation (3 normes)}{9}{10} \\
  \ganttbar[bar/.style={fill=axeC!45,draw=axeC}]
    {J10 — Rapport final \& publication}{10}{11} \\
  \ganttbar[bar/.style={fill=axeC!45,draw=axeC}]
    {J11 — Soutenance finale}{11}{12} \\
  \ganttmilestone[milestone/.style={fill=orange,draw=orange!70!black}]
    {Livrable final}{12}

\end{ganttchart}
\end{center}

\vspace{0.6em}
{\small\textit{%
\textcolor{axeA}{\rule{0.8cm}{5pt}}~Phase 1 / Axe A \quad
\textcolor{axeB}{\rule{0.8cm}{5pt}}~Phase 2 / Axe B \quad
\textcolor{axeC}{\rule{0.8cm}{5pt}}~Phase 3 / Axe C \quad
\textcolor{orange}{\rule{0.3cm}{5pt}\hspace{1pt}$\blacklozenge$}~Jalon de livraison \quad
\textcolor{red!70}{\textbf{- - -}}~Aujourd'hui (Juin 2026).
Ce diagramme sera mis à jour régulièrement.}}

\end{landscape}

\subsection{Méthode de suivi}

Le projet adopte une approche itérative organisée en lots fonctionnels (axes A, B, C),
permettant de structurer les livraisons et leur validation. Le suivi sera assuré via un
dépôt Git, avec un tableau de bord de type Kanban,
mis à jour à chaque feedback encadrants. Les éventuels retards ou changements
technologiques majeurs feront l'objet d'une discussion avec les encadrants pour valider
la faisabilité d'une réorganisation.

\newpage
% ══════════════════════════════════════════════════════════════════════════════
\section{Références bibliographiques}
% ══════════════════════════════════════════════════════════════════════════════
\nocite{*}  % include ALL bib entries in the bibliography
\bibliographystyle{plain}   % ou alpha, ieeetr, abbrv...
\bibliography{bibliography}   % file .bib



\newpage
% ══════════════════════════════════════════════════════════════════════════════
\section*{Historique des versions}
\addcontentsline{toc}{section}{Historique des versions}
% ══════════════════════════════════════════════════════════════════════════════

\begin{center}
\begin{tabularx}{\linewidth}{|C{1.5cm}|C{2.5cm}|>{\raggedright\arraybackslash}X|}
  \hline
  \rowcolor{primary!15}
  \textbf{Version} & \textbf{Date} & \textbf{Description des modifications} \\
  \hline
  1.0 & Juin 2026 & Version initiale du cahier des charges, établie à partir de la présentation intermédiaire. \\
  \hline
  & & \\[0.8em]
  \hline
  & & \\[0.8em]
  \hline
\end{tabularx}
\end{center}

\end{document}